\documentclass[12pt]{article}
\usepackage[T1]{fontenc}
\usepackage[utf8]{inputenc}
\usepackage{lmodern}
\usepackage{textcomp}
\usepackage{enumitem}
\usepackage{geometry}
\geometry{margin=1in}
\begin{document}
\begin{center}
Answer Key - Sociology - Graduate Level Assessment 10 \
\small (Answers are ordered exactly as in the question paper.)
\end{center}

\section*{Multiple Choice}
\begin{enumerate}[label=\arabic*.]
\item \textbf{Answer:} C
\\ \textit{Options:} A) ethnic communities, based on shared identity and experiences of discrimination; B) gay villages, which are formed in certain parts of large cities; C) sociological communities, formed by unpopular lecturers; D) virtual communities that exist only in cyberspace
\\ \textit{Note:} Sociological communities, formed by unpopular lecturers, are not considered a new form of community because they are defined by traditional academic structures rather than emerging social dynamics or innovative forms of collective identity.
\item \textbf{Answer:} D
\\ \textit{Options:} A) the subdivision of labour into small tasks; B) the measurement and specification of work tasks; C) motivation and rewards for productivity; D) all of the above
\\ \textit{Note:} The correct answer is "all of the above" because scientific management encompasses various principles and practices, including standardization of work processes, time studies for efficiency, and the systematic training of workers, all aimed at optimizing productivity in the workplace.
\item \textbf{Answer:} C
\\ \textit{Options:} A) fixed-choice questions; B) short questions; C) leading questions; D) funnelled questions
\\ \textit{Note:} Leading questions can bias respondents' answers by suggesting a particular response or influencing their thoughts, thereby compromising the objectivity and validity of the survey results.
\item \textbf{Answer:} C
\\ \textit{Options:} A) a shared sense of identity and belonging together; B) common activities involving all-round relationships; C) a fixed geographical location; D) collective action based on common interests
\\ \textit{Note:} Fulcher \& Scott argue that community can be defined by social relationships and shared values rather than solely by a fixed geographical location, which is why this option is identified as not a criterion of community.
\item \textbf{Answer:} B
\\ \textit{Options:} A) probability sampling; B) non-probability sampling; C) cluster sampling; D) using the Christmas vacation constructively
\\ \textit{Note:} 'Snowballing' is considered non-probability sampling because it relies on referrals from initial participants to recruit additional subjects, rather than using random selection methods.
\end{enumerate}

\section*{True / False}
\begin{enumerate}[label=\arabic*.]
\setcounter{enumi}{5}
\item \textbf{Answer:} True
\\ \textit{Options:} A) True; B) False
\\ \textit{Note:} Goldthorpe's definition of the 'service class' accurately characterizes this group as comprising individuals in non-manual occupations who perform roles that involve exercising authority and decision-making on behalf of state institutions, distinguishing them from other social classes.
\item \textbf{Answer:} False
\\ \textit{Options:} A) True; B) False
\\ \textit{Note:} Traditional working class identity was primarily characterized by solidarity and distinct cultural values, rather than a homogeneous aspiration to move into the middle class, which often overlooked the complexities of their social and economic realities.
\item \textbf{Answer:} True
\\ \textit{Options:} A) True; B) False
\\ \textit{Note:} Policymakers in the post-Civil War era held the belief that former slaves, liberated from bondage, would successfully integrate into urban society, adopting its cultural norms and participating in economic and social life as part of the broader American melting pot.
\item \textbf{Answer:} True
\\ \textit{Options:} A) True; B) False
\\ \textit{Note:} Decarceration refers to the process of reducing the number of individuals held in prisons or jails, and it includes promoting community-based alternatives such as diversion programs, restorative justice, and mental health services, which aim to address the root causes of criminal behavior while minimizing reliance on incarceration.
\item \textbf{Answer:} True
\\ \textit{Options:} A) True; B) False
\\ \textit{Note:} Modernization theory asserts that societies progress through a linear process of development, transitioning from agrarian, traditional forms of social organization to advanced, industrialized structures characterized by increased technology, urbanization, and complex economic systems.
\end{enumerate}

\section*{Long Form Response}
\begin{enumerate}[label=\arabic*.]
\setcounter{enumi}{10}
\item \textbf{Answer:} Globalization significantly alters the power dynamics of nation-states by introducing transnational forces that challenge their traditional authority. The proliferation of multinational corporations, international organizations, and global communication networks often undermines state sovereignty, as these entities can exert influence that surpasses national borders. Additionally, the rise of global issues such as climate change and migration necessitates collaborative approaches that dilute unilateral state power. Consequently, while nation-states still maintain considerable authority, their ability to act independently is increasingly constrained by the interconnectedness fostered by globalization, leading to a reconfiguration of power that prioritizes global governance over national interests.
\\ \textit{Note:} correct because it identifies how globalization introduces transnational influences and global challenges that erode the traditional sovereignty and authority of nation-states, necessitating cooperation that limits their autonomy.
\item \textbf{Answer:} Blauner's (1964) analysis of alienation in the workplace highlights how the structure of car manufacturing in assembly plants exacerbates feelings of disconnection among employees. Key factors contributing to this alienation include the repetitive nature of assembly line work, which reduces the opportunity for skill development and personal investment in the product, and the lack of autonomy, as workers follow strict protocols and face constant supervision. This mechanization of labor leads to a diminished sense of individuality and creativity, as employees become mere cogs in a vast industrial machine. Furthermore, the hierarchical organization within these plants often results in limited communication and recognition, further alienating workers from their contributions and reducing job satisfaction. Ultimately, Blauner argues that such conditions foster a profound sense of estrangement from both the work process and the final product, perpetuating a cycle of dissatisfaction and disengagement in the automotive industry.
\\ \textit{Note:} The answer correctly identifies and elaborates on Blauner's key concepts of alienation by emphasizing how the repetitive tasks and lack of autonomy in car manufacturing assembly plants contribute to employee disconnection and dissatisfaction.
\end{enumerate}

\vfill
\noindent\textit{End of Answer Key}
\end{document}